\documentclass[11pt]{article}
% ===========================================================================
%  Shared preamble for the AP Calculus AB exam-preparation worksheets.
%  Input right after \documentclass. Files beginning with "_" are not compiled
%  on their own. Navy and gold, to match the course's lesson pages.
% ===========================================================================
\usepackage[letterpaper,top=2.3cm,bottom=2.2cm,left=1.7cm,right=1.7cm,headheight=28pt,footskip=30pt]{geometry}
\usepackage{amsmath,amssymb}
\usepackage[T1]{fontenc}
\usepackage[utf8]{inputenc}
\usepackage{lmodern}
\usepackage{xcolor}
\usepackage[most]{tcolorbox}
\usepackage{enumitem}
\usepackage{fancyhdr}
\usepackage{array}
\usepackage{tikz}
\usepackage{pgfplots}
\pgfplotsset{compat=1.18}

\definecolor{navy}{HTML}{161E54}
\definecolor{ink}{HTML}{27357E}
\definecolor{gold}{HTML}{B98A00}
\definecolor{goldbg}{HTML}{FFFDF3}
\definecolor{inkbg}{HTML}{F6F7FD}
\definecolor{rule}{HTML}{C9CFF0}
\definecolor{okgreen}{HTML}{1B7A44}
\definecolor{okbg}{HTML}{F1F8F3}

\pagestyle{fancy}
\fancyhf{}
\fancyhead[L]{\footnotesize Name: \underline{\hspace{4cm}}}
\fancyhead[C]{\textbf{\color{navy}Integration Academy}\\[-2pt]{\scriptsize Calculus AB \textperiodcentered{} AP\textsuperscript{\textregistered} Exam Preparation}}
\fancyfoot[L]{\tiny\color{black!55}\parbox[t]{0.78\textwidth}{AP\textsuperscript{\textregistered} is a registered trademark of the College Board, which was not involved in the production of, and does not endorse, this material. \textcopyright{} Integration Academy.}}
\fancyfoot[R]{\scriptsize\thepage}
\renewcommand{\headrulewidth}{1.2pt}
\renewcommand{\headrule}{\hbox to\headwidth{\color{gold}\leaders\hrule height \headrulewidth\hfill}}
\renewcommand{\footrulewidth}{0pt}
\setlength{\parindent}{0pt}
\setlength{\parskip}{4pt}

% Title block: \wstitle{code}{title}{unit line}
\newcommand{\wstitle}[3]{\begin{center}{\small\bfseries\color{gold}WORKSHEET #1}\\[3pt]{\LARGE\bfseries\color{navy}#2}\\[3pt]{\small\color{black!60}#3}\end{center}\vspace{2pt}}
% Section heading
\newcommand{\sech}[1]{\par\vspace{10pt}{\large\bfseries\color{navy}#1}\par\nobreak\vspace{2pt}{\color{rule}\hrule height0.8pt}\vspace{7pt}}
% A level tag printed beside a group of questions
\newcommand{\level}[1]{\par\vspace{6pt}{\footnotesize\bfseries\color{gold}\MakeUppercase{#1}}\par\nobreak\vspace{2pt}}

\newtcolorbox{ideabox}{colframe=ink,colback=inkbg,coltitle=white,colbacktitle=ink,fonttitle=\bfseries\small,title={Key ideas},arc=2pt,boxrule=0.8pt,left=8pt,right=8pt}
\newtcolorbox{example}[1]{colframe=rule,colback=white,coltitle=navy,colbacktitle=inkbg,fonttitle=\bfseries\small,title={#1},arc=2pt,boxrule=0.8pt,borderline west={3pt}{0pt}{ink},left=10pt,right=8pt}
\newtcolorbox{frbox}[1]{colframe=gold,colback=goldbg,coltitle=white,colbacktitle=gold,fonttitle=\bfseries\small,title={#1},arc=2pt,boxrule=0.8pt,breakable,left=8pt,right=8pt}
\newtcolorbox{solbox}{colframe=okgreen,colback=okbg,coltitle=white,colbacktitle=okgreen,fonttitle=\bfseries\small,title={Worked solutions},arc=2pt,boxrule=0.8pt,breakable,left=8pt,right=8pt}
\newtcolorbox{scorebox}[1]{colframe=navy,colback=white,coltitle=white,colbacktitle=navy,fonttitle=\bfseries\small,title={#1},arc=2pt,boxrule=0.8pt,breakable,left=8pt,right=8pt}

\newcommand{\soln}{\par\smallskip\textbf{\color{ink}Solution.}\ }
\newcommand{\result}[1]{\par\smallskip\textbf{#1}}
% blank working space of a given height
\newcommand{\work}[1]{\par\vspace{#1}}
% a point value in the margin of a scoring guide
\newcommand{\pts}[1]{\hfill{\footnotesize\bfseries\color{gold}[#1]}}

% Coordinate grid for graph questions: \begin{apgraph}{xmin}{xmax}{ymin}{ymax} ... \end{apgraph}
\newenvironment{apgraph}[4]{\begin{center}\begin{tikzpicture}\begin{axis}[
  axis lines=middle, xmin=#1, xmax=#2, ymin=#3, ymax=#4,
  grid=both, grid style={black!12}, xtick distance=1, ytick distance=1,
  ticklabel style={font=\scriptsize}, xlabel={$x$}, ylabel={$y$},
  width=10cm, height=7.2cm, samples=80, clip=false,
  every axis plot/.append style={line width=1.1pt, color=ink}]}{\end{axis}\end{tikzpicture}\end{center}}
\newcommand{\closeddot}[2]{\addplot[only marks, mark=*, mark size=2.2pt, color=ink] coordinates {(#1,#2)};}
\newcommand{\opendot}[2]{\addplot[only marks, mark=*, mark size=2.2pt, color=ink, mark options={fill=white}] coordinates {(#1,#2)};}

% Multiple-choice options: four across, or stacked when the options are long.
\newcommand{\opts}[4]{\par\vspace{3pt}\begin{tabular}{@{}p{0.235\linewidth}p{0.235\linewidth}p{0.235\linewidth}p{0.235\linewidth}@{}}(A)\ #1 & (B)\ #2 & (C)\ #3 & (D)\ #4\end{tabular}\par\vspace{9pt}}
\newcommand{\optsv}[4]{\par\vspace{3pt}\begin{tabular}{@{}l@{\ \ }p{0.9\linewidth}@{}}(A) & #1\\(B) & #2\\(C) & #3\\(D) & #4\end{tabular}\par\vspace{9pt}}

\fancyhead[R]{\footnotesize\color{black!60}Worksheet 1.8}
\begin{document}

\wstitle{1.8}{The Three-Part Continuity Test}{Unit 1: Limits and Continuity}

\begin{ideabox}
\begin{itemize}[leftmargin=*,itemsep=2pt,topsep=2pt]
  \item $f$ is continuous at $c$ when $f(c)$ is defined, $\displaystyle\lim_{x\to c}f(x)$ exists, and the two are equal.
  \item For a piecewise function the justification is one line: $\displaystyle\lim_{x\to c^-}f(x)=\displaystyle\lim_{x\to c^+}f(x)=f(c)$.
  \item Polynomial, rational, root, trigonometric, exponential and logarithmic functions are continuous on their domains.
  \item To make pieces join, set the limits equal at each boundary: one equation per boundary.
\end{itemize}
\end{ideabox}

\sech{Worked examples}
\begin{example}{Example 1 \textperiodcentered{} Verifying continuity}
Is $f(x)=\begin{cases}x^3,&x\le1\\3x-2,&x>1\end{cases}$ continuous at $x=1$?\soln $\displaystyle\lim_{x\to 1^-}f(x)=1$, $\displaystyle\lim_{x\to 1^+}f(x)=3-2=1$, and $f(1)=1$.\result{Yes: the three values are equal.}
\end{example}

\begin{example}{Example 2 \textperiodcentered{} One constant}
Find $k$ so that $f(x)=\begin{cases}kx-3,&x<2\\x^2+k,&x\ge2\end{cases}$ is continuous at $x=2$.\soln $2k-3=4+k$.\result{$k=7$. Check: both sides give 11.}
\end{example}

\begin{example}{Example 3 \textperiodcentered{} A function value and two constants}
Find $a$ and $b$ so that $f(x)=\begin{cases}ax+2,&x<1\\5,&x=1\\bx^2,&x>1\end{cases}$ is continuous at $x=1$.\soln Left limit: $a+2$. Right limit: $b$. Both must equal $f(1)=5$.\result{$a=3$ and $b=5$.}
\end{example}

\newpage
\sech{Practice}

\level{Core}
\begin{enumerate}[leftmargin=*,itemsep=2pt]
  \item Is $f(x)=\begin{cases}4-x,&x<3\\x-2,&x\ge3\end{cases}$ continuous at $x=3$? Justify.\work{2.3cm}
  \item On what set is $f(x)=\dfrac{1}{x^2-9}$ continuous?\work{2.3cm}
  \item On what interval is $f(x)=\sqrt{x-2}$ continuous?\work{2.3cm}
  \item What value should be assigned to $f(6)$ to make $f(x)=\dfrac{x^2-36}{x-6}$ continuous at $x=6$?\work{2.3cm}
\end{enumerate}

\level{Exam level}
\begin{enumerate}[leftmargin=*,itemsep=2pt,start=5]
  \item Find $k$ so that $f(x)=\begin{cases}x^2-k,&x\le2\\kx+1,&x>2\end{cases}$ is continuous at $x=2$.\work{2.8cm}
  \item Find $k$ so that $f(x)=\begin{cases}\dfrac{\sin(4x)}{x},&x\neq0\\k,&x=0\end{cases}$ is continuous at $x=0$.\work{2.8cm}
  \item Find $a$ and $b$ so that $f(x)=\begin{cases}x+a,&x<0\\x^2+3,&0\le x\le2\\bx-1,&x>2\end{cases}$ is continuous for all real $x$.\work{2.8cm}
  \item On what interval is $f(x)=\ln(5-x)$ continuous?\work{2.8cm}
\end{enumerate}

\begin{frbox}{Question 9 \textperiodcentered{} Free-response practice}
Let $f(x)=\begin{cases}x^2+1,&x<2\\5,&x=2\\7-x,&x>2.\end{cases}$
\begin{enumerate}[label=(\alph*),leftmargin=*,itemsep=2pt]
  \item Find $\displaystyle\lim_{x\to 2^-}f(x)$ and $\displaystyle\lim_{x\to 2^+}f(x)$.\work{2.6cm}
  \item Is $f$ continuous at $x=2$? Justify your answer.\work{2.6cm}
  \item Let $h(x)=f(x)$ for $x\neq2$ and $h(2)=4$. Is $h$ continuous at $x=2$? If not, name the type of discontinuity.\work{2.6cm}
\end{enumerate}
\end{frbox}

\newpage
\begin{solbox}
\begin{enumerate}[leftmargin=*,itemsep=4pt]
  \item $\displaystyle\lim_{x\to 3^-}f(x)=1$, $\displaystyle\lim_{x\to 3^+}f(x)=1$, $f(3)=1$. All equal, so yes.
  \item All real numbers except $x=3$ and $x=-3$.
  \item $[2,\infty)$.
  \item $12$, the limit of $x+6$.
  \item $4-k=2k+1$ gives $k=1$.
  \item $\displaystyle\lim_{x\to 0}\dfrac{\sin4x}{x}=4$, so $k=4$.
  \item At 0: $a=3$. At 2: $7=2b-1$, so $b=4$.
  \item $5-x>0$: the interval $(-\infty,5)$.
  \item \begin{enumerate}[label=(\alph*),leftmargin=*,itemsep=3pt]
    \item $\displaystyle\lim_{x\to 2^-}f(x)=4+1=5$ and $\displaystyle\lim_{x\to 2^+}f(x)=7-2=5$.
    \item Yes. $\displaystyle\lim_{x\to 2^-}f(x)=\displaystyle\lim_{x\to 2^+}f(x)=f(2)=5$.
    \item No. $\displaystyle\lim_{x\to 2}h(x)=5$ but $h(2)=4$: a removable discontinuity.
  \end{enumerate}
\end{enumerate}
\end{solbox}

\end{document}
